% Sin número: se conserva la numeración de las secciones del documento.
\section*{Concepto de Blueprint}
\addcontentsline{toc}{section}{Concepto de Blueprint}

Un Blueprint es un plano o guía que define y orienta la transformación de una organización hacia un
estado futuro deseado (\textit{target state}) \parencite{moyle2020practical}. Reúne y organiza los
elementos necesarios para diseñar una mejora y preparar su implementación, de modo que las decisiones
de diseño respondan a los objetivos del negocio y a los aspectos tecnológicos y de seguridad
involucrados.

Para elaborarlo se necesita información previa sobre la organización, sus procesos, objetivos,
necesidades y riesgos. TOGAF ofrece una perspectiva de la arquitectura empresarial y de la forma en que
se estructuran sus procesos y componentes; SABSA permite analizar los requerimientos y riesgos de
seguridad y relacionarlos con las capas del negocio \parencite{moyle2020practical}. Esa información
sirve de base para un Blueprint alineado con las necesidades y los objetivos de la organización.

Un Blueprint fija los principios de arquitectura que orientan las decisiones de diseño, delimita el
alcance, describe el estado futuro y alinea la arquitectura con los objetivos del negocio. También
sirve de guía para la implementación, facilita la comunicación y validación con las partes interesadas
y deja una referencia para hacer seguimiento y medir el avance.

\textcite{moyle2020practical} comparan el Blueprint con el plano de una casa. Antes de
construir hay que conocer qué quiere el cliente, cuál es el presupuesto, qué características tiene el
terreno, cuánto tiempo hay y qué condiciones deben cumplirse; con esa información se elabora el plano
que guía a los especialistas. Un Blueprint hace lo mismo con la organización: convierte una visión
inicial en un diseño que puede comprenderse, validarse e implementarse. Por eso reúne modelos,
documentos y artefactos que fijan una visión clara del estado futuro y una guía para alcanzarlo de
forma segura, consistente y alineada con los objetivos del negocio \parencite{moyle2020practical}.
